\section{조건부 보안 증명의 세부사항}
\label{app:security-details}

이 부록은 VI절의 세 정리에 사용한 실험과 암호 가정을 명시한다.
모든 공격자와 실행·객체·증명·질의 수는 보안 매개변수 $\lambda$의
다항식으로 제한한다. 확률은 명시한 전체 실행에 관한 것이다.
원장의 권한·인증·확정·현재 상태 집행과 신뢰된 감사자의 동작은
모델 전제이며 아래 암호 가정으로 새로 증명하지 않는다.

\subsection{정확성 증명}
\label{app:correctness-details}

정확성 게임은 정직한 생성에 사용한 증인을 보관한다.
거래는 전체 증명 관계와 실행 조건, 출력 신규성과 트리 용량을
만족하며 그 조건이 실행까지 유지되어야 한다.
공개 DPP 검증은 발급이 포함된 확정 이력과 그 발급에 맞는 정규 문서를 쓴다.
추적은 유효 시작점·확정 $H_t$·유효 감사키를 사용하고,
필요한 유한 기록과 종료하는 인증 조회가 있어야 한다.
고정 $H_t$의 조회에는 그 이후 원장 전체의 무경합을 요구하지 않는다.

차단 승인이 있고 정확한 전체 $T$가 비어 있지 않다면,
$H_t$부터 $H_r$까지 대상 소비와 차단 해제가 없어야 한다.
Active 대상의 합법적 동결, 모든 제출 결과·조회,
사용한 관찰을 반영하는 같은 확정 이력의 $H_r$을 정한 유한 한도
안에서 확보해야 한다. $T=\varnothing$에는 이 추가 조건 없이
빈 대상 분기를 적용한다.

위 조건이 성립하는 요청에서 기대한 결과를 내지 못하면
$\mathsf{Fail}_{\rm corr}$를 기록한다. 조건이 깨져 정상 거부하거나
\Alg{INCOMPLETE}인 실행은 이 사건이 아니다.
증명 완전성 실패 밖에서는 정직 증명이 수락되고 원자적 상태 갱신이
정해진 결과를 기록한다. 복호화 실패 밖에서 실제 부모·미래 식별자를
얻고 관련 해시 충돌 밖에서 다른 객체의 식별자는 구별된다.
출력 참조의 신규성만으로 식별자 충돌을 배제할 수는 없다.

부록~\ref{app:trace-details}의 경로를 정직 생성 목록에 적용하면
그래프·미소비 집합·종단이 정확하다.
증인을 이미 알고 있으므로 정확성에는 악의적 증인의 추출이 필요 없다.
동결의 추가 조건은 전체 $T$와 공통 시점의 완료 검사를 만족시킨다.
따라서 실패는 완전성·복호화·관련 충돌의 합집합에 포함되고,
정리 1의 전체 실행 상계가 따른다.

\subsection{건전성 증명}
\label{app:soundness-details}

보안 강도의 기준값인 보안 매개변수 $\lambda$를 두고,
그 값에 대한 다항 길이 실행을 고려한다.
공격자는 정상 거래를 관찰하거나 만들게 하고,
이전 결과를 본 뒤 다음 행동을 고를 수 있다.
임의의 공개 입력과 증명 $(x,\pi)$도 제출한다.
확정된 과거 이력과 문서·Claim에 대해 공개 검증 절차인
\Alg{CheckDPP}도 질의한다. 원장은 고정된 전체 관계의 검증키로 증명을 확인하고 현재
권한·정책·미소비·활성·기한·출력 신규성을 검사한 뒤 원자적으로
갱신한다. 수락 이력을 처음부터 의미적으로 유효하다고 가정하지 않는다.

나중에 제출된 값이 실제 생성값과 같은지 비교하려면
처음 생성된 값을 기준으로 남겨야 한다.
이를 위해 분석용 생성 목록 $\mathcal C$에 객체가 처음 수락됐을
때의 전체 개방값과 타입별 미래 식별자를 고정한다. 정직 생성에는 원래 증인,
공격자 수락에는 추출 증인을 사용하며 식별자에서 생성 객체로 가는
역대응도 고정한다. 이는 구현에 추가하는 공개 목록이 아니다.
실패 사건 $\mathsf{BadSem}$은 수락 입력이 먼저 생성되지 않았거나,
생성 개방값과 다른 값을 사용하거나, 그 실제 값에 대한 소유 비밀·
정수 회계·정책 관계 또는 현재 상태 검사가 거짓인 경우다.
반복 계상·소비, 잘못된 상태 갱신, 실제 부모·미래 식별자와 다른
감사 복호화도 포함한다. 공개 DPP는 인증된 발급 기록과 입력
문서 해시·정책·논스의 불일치를 실패로 정의한다.
추출 실패나 모호한 역대응에서는 분석 게임이 실패를 기록하고
종료한다. 그 밖에서만 실제 생성 목록과 의미 위반을 비교한다.

공격자는 수락 이력의 유효 시작점·확정 시점·방향을 골라 추적을
실행하게 할 수 있다. 차단 승인이 있는 후속 추적에는 신뢰된 상태
관리자가 \Alg{AuditAndFreeze}를 수행한다. 공격자는 허용된 전이·
합법적 상태 변경·자료 오류를 스케줄하며 감사 이후의 전이까지
전체 실행 예산에 포함한다. 인증 조회는 실제 값 또는 오류만
제공하고, 임의 권한 변경이나 새로운 키 공개 질의는 주지 않는다.

통합 게임 $\mathsf{Exp}_{\rm sound}$는 위 전이·DPP 의미 위반,
잘못된 추적 성공, 또는 완료 반환과 완료 술어/차단 유지 구간의
위반이 함께 생기면 1을 반환한다. 자료 부재의 명시적 실패나 전체
대상을 보존한 \Alg{INCOMPLETE}는 공격 성공이 아니다. 아래 나쁜
사건에서는 1로 종료하고 그 밖에서만 생성 목록과 결과를 비교한다.

증인 추출·결합·포함 관계·해시·복호화에 다음 계약을 적용한다.
\begin{enumerate}
 \item 정직한 증명의 관측·재전송과 다수 등록 관계를 포함하는
 적응적 실행의 수락 증인을 공동 추출하는 PPT 알고리즘이 있다.
 추출은 분석용 권한으로 증명의 근거 입력을 얻는다는 뜻이다.
 일반 검증자가 비공개 입력을 볼 수 있다는 뜻은 아니다.
 비공개 추출 보조키(trapdoor) 또는 명시된 공격자 접근 권한을 가지며
 이를 관측자에게 주지 않는다. 공개 실행 기록의 분포는 실제 실행과
 같거나 통계적으로 가까우며 그 오차도 전체 실패 상계
 $\epsilon_X$에 포함한다. 추출 보조정보는 독립 원시기술의
 도전 실험의 비밀을 주지 않아야 한다. 정직한 동일 증명
 재전송에는 원래 증인을 쓴다. 단일 증명의 지식 건전성만으로
 이 계약을 추정하지 않는다.
 \item 정규 전체 개방값의 계산적 결합성, 수락된 루트의 포함
 관계의 건전성, 관련 객체·식별자·Claim 해시의 충돌 저항성에
 대한 전체 실행 상계를 각각 $\epsilon_C,\epsilon_M,\epsilon_H$라
 둔다. 유효 비밀 영역에서 $\mathsf{OwnerKey}$는 단사이다.
 즉 서로 다른 유효 소유권 비밀이 같은 소유자 값으로 바뀌지 않는다.
 \item 고정 키·문맥에서 모든 허용 평문과 \emph{증명자가 선택한}
 유효 난수에 대해 복호화가 정확하다. 완전 정확성이면 실패는 0이며,
 완화할 경우 적대적 난수 선택의 실패 상계를 $\epsilon_E$로 둔다.
 정직 난수 평균에서만 정확하다는 조건은 충분하지 않다.
\end{enumerate}
소유권 키의 예로 소수 위수 $\ell$의 생성점 $G$에 대해 $sG$와
정규 영역 $1\leq s<\ell$를 사용할 수 있다. 단사성과 비밀의
일방향성은 다른 성질이다. 원장·인증·권한·정책 구성의 정확성은
모델 전제이며 위 암호 가정으로 새로 증명하지 않는다.

\emph{증명: 전이와 감사 결합.}
추출·결합·포함 관계·관련 충돌·복호화 계약의 실패 합집합을
$\mathsf{Bad}$라 두면 $\Pr[\mathsf{Bad}]\leq\epsilon_B$다.
이 중 어느 실패도 일어나지 않은 실행을 생각한다.
앞서 수락된 기록이 올바르면 다음 거래도 그 기록을
올바르게 사용해야 함을 수락 순서대로 보인다.
추출 입력의 포함 증인은 이전 실제 등록 객체를 가리킨다.
생성 목록의 개방값과 다르면 계산적 결합성의 명시적 충돌 쌍이다.
따라서 추출 증인의 타입·정수·소유권·정책 관계가 바로 최초
생성 값에 적용된다. 계산적 결합성을 모든 수학적 증인의
유일성으로 해석하는 것은 아니다.

Note의 생성·소비 개방값은 같은 $cm,owner$를 포함한다.
단사성으로 소유 비밀도 같아 $nf$가 같다. Voucher는 결합된 같은
개방 난수에서 $rvnf$를 얻으므로 인수·회수도 동일하다.
다른 같은 타입 객체의 식별자 중복은 서로 다른 참조가 포함된
정규 해시 입력의 충돌이다. 이에 따라 타입별 현재 소비·상태
검사가 바로 실제 객체에 적용된다. 중복 입력 검사와 원자적
소비 기록은 반복 계상과 경합 재소비를 거부한다.
과거 포함 관계의 루트는 현재 상태 검사를 대체하지 않는다.

감사 관계는 이 실제 입력·출력의 부모·미래 식별자를 같은 증인으로
암호화한다. 고정 키·문맥·슬롯과 적대적 난수의 복호화 정확성으로
감사자는 그 평문을 얻는다. 다른 문서 해시·정책·논스를 기존
등록 Claim에 붙여 유효 판정을 받으면 Claim 결합 해시의 충돌이다.
따라서 실패 합집합 밖에서는 $\mathsf{BadSem}$이 없고,
같은 나쁜 사건 밖에서 추적·차단을 분석할 수 있다.

\paragraph{추적과 전체 대상의 차단.}
실제 그래프는 위 생성 목록으로 부록~\ref{app:trace-details}처럼 정의한다.
공격자는 수락 이력 뒤 확정 시점·유효 시작점·방향과 자료 오류를 고른다.
인증 조회는 그 시점의 실제 값 또는 오류만 제공하고 감사자에게
유효한 키를 준다. 나쁜 사건에서는 실패를 기록하고,
그 밖에서는 성공 반환이 실제 그래프·미소비 집합·종단과 다른지 비교한다.

부모 복호화는 실제 입력을, 미래 식별자는 실제 소비 이벤트를 가리킨다.
시작점에서 한 단계 확장이 맞으면 다음 단계도 맞으므로,
경로 길이에 대한 귀납으로 추적 결과의 일치를 얻는다.
재방문에서도 새 간선을 남기므로 분기·합류에서 경로를 잃지 않는다.
필요한 자료가 가용하고 조회가 종료하면 유한 그래프의 추적도 끝난다.
항상 실패하는 절차를 정확한 감사로 취급하지 않는다.

감사 자료에 출력과 무관한 식별자를 넣거나 인증되지 않은 0 응답을
허용하면 다른 객체의 소비 여부를 조회하거나 소비된 객체를
미소비로 오인할 수 있다. 일부 소비가 현재 상태 검사를 빠뜨리면
동결 뒤에도 그 경로를 사용할 수 있다. 이는 가정 생략의
모델 반례이며 특정 구현의 취약점 보고는 아니다.

완료 게임에서는 공격자가 합법적 거래·상태 변경과 자료 오류를
스케줄하고 신뢰된 상태 관리자가 실제 \Alg{AuditAndFreeze}를 실행한다.
감사 이후 이벤트도 전체 실행 예산에 포함한다.
공통 $H_r$은 같은 확정 이력에서 사용한 관찰·영수증을 반영한다.
정확한 전체 $T$에 식~\eqref{eq:frontier-cut}가 성립하면,
$H_t$ 뒤 최초의 후속 소비는 $T$의 객체를 써야 한다.
그러면 그 불변 소비 기록이 $H_r$에서 0일 수 없어 모순이다.
$H_r$ 뒤에도 차단이 유지되면 현재 상태 검사가 첫 소비를 거부한다.

완료를 반환했으면서 완료 술어가 거짓이거나 유지 구간의 차단 함의가
깨지는 사건을 $\mathsf{BadComplete}$라 하자.
위 결정적 함의로 $\Pr[\mathsf{BadComplete}]\leq\epsilon_B$를 얻는다.
상계는 나쁜 사건을 포함한 전체 실행에 관한 것이며
완료에 조건을 건 확률로 바꾸지 않는다.
전이·추적·완료의 실패는 같은 나쁜 사건에 포함되므로
정리 2의 $\epsilon_B$를 세 번 더하지 않는다.

전체 차단은 본 절차의 충분조건이다.
대상 종료로 경로가 소진됐거나 다른 실행이 새 대상을 막았더라도
원래 $T$의 완료 판정과 같지는 않다.
모든 차단이 되돌릴 수 없는 사용 금지이고 인증된 관찰을 이후
공통 상태로 정당하게 확장할 수 있다면 다른 확인 절차도 가능하다.
본 결과는 모든 프로토콜의 재조회 필요성, 영구 차단·Claim 취소·
물리 회수나 모든 대응 방식의 종료 불가능성을 주장하지 않는다.

\subsection{프라이버시 증명}
\label{app:privacy-details}

\paragraph{두 키 암호화의 기밀성 보조정리.}
\label{sec:wrapper-proof}

암호문이 내용을 숨기는지 보려면 두 원문 중 어느 것을
암호화했는지 맞히게 할 수 있다.
원문만 바꾼 두 실험에서 공격자가 1을 출력할 확률이
얼마나 달라지는지를 \emph{구별 이득}이라고 한다.
차이가 작을수록 암호문만으로 두 경우를 구분하기 어렵다. 독립적으로 생성한 A/B 공개키와 설정 전에 고정한 $K-1$개
이하의 위원회 조각을 공개한다. 공격자는 같은 길이 감사 평문 쌍을
적응적으로 선택하고 한쪽의 정직한 암호문 및 $W_B$를 받는다.
$SK_B$, 복호화·키 복구와 도전 암호문 $W_B$의 부분 복호화
질의는 제공하지 않는다. 인코딩 길이는 고정되고, 같은 길이의
공개적으로 계산할 수 있는 유효한 대체 평문이 있다. 공개 설정 정보와 인증은
각 환원에서 재현 가능하며, 숨은 키 일치에 대한 별도의 비시뮬레이션
증거를 포함하지 않는다. 여러 암호문을 함께 보는 공개키 암호 게임과,
노출할 조각을 설정 전에 고정하는 임계 암호 게임을 사용한다\cite{boneh-shoup-2020}.

\textbf{보조정리 (두 키 암호화).}
TE는 임계 암호화, PKE는 일반 공개키 암호화를 뜻한다.
CPA는 공격자가 비교할 원문을 선택할 수 있는 보안 실험이다.
TE의 해당 CPA 이득을 $\epsilon_A$, 여러 데이터 암호문을 함께
보는 PKE의 CPA 이득을 $\epsilon_D$라 하면 이 게임의 이득은
$2\epsilon_A+\epsilon_D$ 이하이다.

\emph{증명.}
공격자가 보는 자료를 세 단계로 바꾼다.
각 단계의 변화를 구별할 수 있다면 그 단계에 사용한
암호 가정도 깨뜨릴 수 있다.
먼저 $W_B$를 실제 $SK_B$ 대신 대체 평문의 암호화로 바꾼다.
TE 환원은 B 키쌍을 스스로 생성하므로 두 평문을 안다.
다음 데이터 암호문을 반대편 평문으로 바꾼다. PKE 환원은 받은
$PK_B$의 $SK_B$를 몰라도 대체 평문을 담은 키 보호 암호문을 생성할 수 있다.
마지막으로 실제 키를 담은 암호문을 복원한다. 세 변화의 삼각 부등식이
결론이다. 두 독립 키를 쓰므로 이 구조 자체에 자기 키 암호화
가정을 추가할 필요는 없다. 이 결과는 다른 공개 필드의 비밀성이나
외부 키 전달 서비스의 인증·강건성을 증명하지 않는다. $\square$

\paragraph{원장 구별 게임과 충분조건.}
\label{sec:passive-privacy}

전체 원장의 비밀성에는 커미트먼트·소비 식별자·증명의 공동
분포도 필요하다. 특히 이전 증표는 같은 개방 난수를 커미트먼트와
해결 비밀값에 사용하므로 개별 은닉성과 의사난수성만으로
그 합성의 보안을 결론 내릴 수 없다. 다음은 이 공동 분포에 대한
명시적 충분조건이며, 구체 해시·회로의 충족 여부는 별도 검증 대상이다.

여기서는 해시를 이상적인 무작위 함수로 다루는
\emph{무작위 오라클 모형}을 사용한다.
처음 보는 입력에는 무작위 값을 정하고, 같은 입력에는
항상 같은 값을 돌려주는 모형이다.
회계·이전 증표의 커미트먼트, 해결 비밀값과 두 소비 식별자의
해시 영역을 서로 독립된 고전적 무작위 오라클로 모형화한다.
이는 실제 해시 함수의 동작을 그대로 증명했다는 뜻은 아니다.

출력 영역의 크기는 최소 $2^\lambda$, 개방 난수와 정규 소유 비밀의
독립 균등 영역은 최소 $2^\kappa$이고 $\kappa$는 $\lambda$에
선형으로 증가한다. 소유 비밀 재사용은 허용하되 이 게임에서
소유 비밀·개방값·별도 공개 OwnerKey는 제공하지 않는다.
공개 호출자·서명키 등은 보호 비밀과 독립이다.

영지식 성질(ZK)은 분석용 보조 정보를 가진 알고리즘이
비공개 증인을 몰라도 실제 증명과 구별하기 어려운 자료를
만들 수 있다는 방식으로 분석한다. 이 알고리즘이 시뮬레이터다.
일반 검증자에게 증명을 위조할 권한을 준다는 뜻은 아니다.
두 키 보조정리의 암호 가정과 함께, 이 공동 오라클 모형에서
여러 증명과 보조 정보를 함께 보더라도 안전한 다중 정리 ZK를
가정한다. 시뮬레이터는 모든 정규 문장에서 종료하는
PPT $\mathsf{SimProof}(td,x)$이다.

공격자는 공개 설정을 본 뒤 비교할 두 이력 $h_0,h_1$을 고른다.
예를 들어 숨긴 회계 수치나 실제 입력 연결이 서로 다를 수 있다.
다만 공개하기로 한 정보는 두 이력에서 같아야 한다.
이력은 회계 값·소유자 이름·연결을 지정하지만 비밀·난수를
지정하지 않는다. 이 공개 정보의 범위를 공개 누출 $L$이라고 한다.
$L$은 종류·정책·
적용 범위·타입·입출력 수·순서·동일 이벤트의 출력 관계·호출자·시간·기한·루트의
접두 이력 내 위치·암호문 길이·공개 Claim의 문서 해시와 논스 의미다.
공개 객체와 소비 식별자의 타입별 표지, 생성 이벤트·색인 구조와
동일성 패턴도 포함한다. 비공개인 생성 참조에서 미래 식별자로 가는
대응은 $L$에서 자동 연결하지 않는다. 원시 바이트가 같아야 한다는
요구는 아니다.

선택한 이력은 독립 정직 난수로 실현되며, 우연한 충돌 외에는 그
난수에 따른 선택적 수락·비트 조건·재표본화를 요구하지 않는다.
관측자는 모든 공개 원장·암호문·증명을 보고 공개 오라클 질의를
수행하지만, 참여자 부패·개방값 공개·능동 제출·복호화 또는
보호 관계의 감사·상태 공개 질의는 받지 않는다. 위원회 노출은
설정 전 고정한 $K-1$개 이하로 유지한다.

위 게임에서 전체 증명의 ZK 이득을 $\epsilon_Z$라 하면
정리 3의 식~\eqref{eq:privacy-bound}가 성립한다.
객체 수 $n$, 소유자 수 $u$, 외부·시뮬레이터의 관련 질의 수 $Q$에
대해 $\epsilon_{\rm RO}=O((n+u+Q)^2/2^{\min(\kappa,\lambda)})$의
보수적 상계를 사용할 수 있다.
이 식은 다루는 객체와 질의가 많아질수록, 비밀을 고르는 영역이
작아질수록 구별 가능성의 한도가 커진다는 뜻이다.
구체 구현의 보안 비트나 실행 시간을 측정한 수치는 아니다.

\emph{증명.}
각 이력에서 먼저 참인 공개 문장의 증명을 ZK 시뮬레이터로
바꾼다. 다음 두 키 보조정리의 순서로 키 보호 암호문의 평문과
감사 평문을 대체 평문으로 치환한다. 마지막으로 숨은 커미트먼트와
미래 식별자에 독립 무작위 표지를 배정하되 공개 재사용·동일성은 보존한다.
회계 증표의 커미트먼트 입력에는 숨은 개방 난수가, $nf$ 입력에는
소유 비밀이 필요하다. 같은 비밀을 재사용해도 다른 $cm$는
다른 오라클 입력이다. 이전 증표의 두 영역은 같은 개방 난수에
의존하나 해결 비밀값은 별도 균등 값이다. 그 개방 난수
또는 해결 비밀값을 맞히기 전에는 공동 오라클 표와
독립 표지를 같은 분포로 결합할 수 있다.

비밀을 받지 않은 외부·시뮬레이터의 첫 숨은 입력 질의와 관련
충돌을 나쁜 사건으로 두면 균등 비밀 추측 및 생일 충돌 상계가
$\epsilon_{\rm RO}$를 준다. 그 밖의 루트·색인과 암호문 문맥은
공개 표지의 함수다. 마지막 분포는 동일 $L$과 독립 키·표지로만
정해지므로 양쪽 중간 분포에 삼각 부등식을 적용하면 정리가 따른다.
암호문 치환 후 문장이 거짓이어도 ZK를 다시 적용하지 않는다.
각 환원은 모든 입력에서 종료하는 PPT $\mathsf{SimProof}(td,x)$를 효율적 후처리로
실행할 뿐 중간 원장의 유효성을 주장하지 않는다. $\square$

이 정리는 일반 참여자 프라이버시 전체와 다르다. 소유 비밀을
공개하면 그 비밀로 과거·미래 회계 증표의 식별자를 대입할 수 있고,
이전 증표의 개방값을 아는 송·수신자는 그 해결 식별자를 안다.
부패 참여자의 지식, 적응적 공개, 공개 OwnerKey를 허용하는 확장과
실제 해시·증명계의 공동 모형 적합성에는 별도의 증명이 필요하다.
